\clearpage
\thispagestyle{empty}
\pdfbookmark[0]{Anleitung zur Bestimmung der Schwingungsdauer einer Magnetnadel. Ausgabe und Quellen}{magnetic-needle-swing-sources}
\begin{center}\Large\textbf{Anleitung zur Bestimmung der Schwingungsdauer einer Magnetnadel}\end{center}

Carl Friedrich Gauss’ Anleitung erschien in den \textit{Resultaten aus den Beobachtungen des magnetischen Vereins} für 1837 als Artikel IV. Diese KI-Übertragung in modernes Deutsch mit KI-erstelltem Satz folgt dem vollständigen deutschen Text in den \textit{Werken}, Bd.~V (Göttingen, 1867), mit dem unnummerierten Anfang, S.~375--394 und den fünf historischen Anmerkungen. Eine Schwingung reicht hier von einer äußersten Stellung bis zur nächsten. Die Beobachtungen und Formeln stehen in der historischen Notation und in den historischen Einheiten.

\smallskip\noindent
Zur Terminologie der Durchgangsinstrumente und der Zeitbestimmung siehe den Katalogeintrag \textit{Passageinstrument} des Deutschen Museums, herausgegeben von Julia Bloemer und Benjamin Mirwald, Nr.~701, Fassung vom 24. August 2016, sowie die Glossareinträge \textit{Auge-Ohr-Methode} und \textit{Fadenkreuz}. Zur modernen Terminologie gedämpfter Schwingungen siehe die unten verlinkten Physik-Praktikumsanleitungen des Karlsruher Instituts für Technologie.

\medskip\noindent
\href{https://gdz.sub.uni-goettingen.de/id/PPN236006339}{Gesammelte deutsche Ausgabe (1867)}\par
\noindent\href{https://digital.deutsches-museum.de/projekte/gruendungssammlung/detail/701/}{Katalogeintrag zum Passageinstrument}\par
\noindent\href{https://digital.deutsches-museum.de/projekte/gruendungssammlung/glossar/}{Glossar des Deutschen Museums}\par
\noindent\href{https://labs.physik.kit.edu/prakt-klass-physik-manuals/P1/Resonanz/doc/Hinweise-Schwingung/}{KIT, Praktikum Klassische Physik, Hinweise zu Schwingungen}\par
\noindent\href{https://labs.physik.kit.edu/prakt-nebenfach-manuals/Experiments/Erzwungene_Schwingungen/}{KIT, Nebenfach-Praktikum, Erzwungene Schwingungen}
\clearpage
